\documentclass[12pt]{article}
\usepackage[utf8]{inputenc}
\usepackage[T1]{fontenc}
\usepackage{hyperref}
\title{Xinca Language Resources for K-12 Curriculum Development}
\author{Compiled by Claude Sonnet 5.5 (claude-sonnet-5-5) for Andr\'as Kornai}
\date{}
\begin{document}
\maketitle

\section{Introduction}
Xinca (ISO 639-3 xin) is a dormant language of Guatemala. Britannica describes a Xinkan group of four languages, two of them extinct, and notes that some young people are learning Guazacapán Xinka as a second language. A 2003 decree recognised Xinka as a national language. Documentation is thin. The sources below are mostly historical wordlists and academic work, not school materials.

\section{Primary Educational Resources}
I found no primers, textbooks, or school programmes for Xinca. One Webonary Xinka site exists (see below), but I could not confirm what it contains beyond its introduction.

\section{Linguistic and Reference Materials}
\begin{itemize}
\item Xinka site on Webonary (introduction page); the extent of its dictionary content is unverified. \url{https://www.webonary.work/xinka/overview/introduction/}
\item Endangered Languages Project entry for Xinca, listing historical materials such as an 1868 vocabulary compiled by C.\ Hermann Berendt. \url{https://endangeredlanguages.com/node/112966}
\item Frauke Sachse, publications, including work on the colonial \emph{Arte de la lengua szinca} and fieldwork in Guazac\'apan. \url{https://doaks.academia.edu/FraukeSachse/Books}
\item Leiden University repository item on Xinkan. \url{https://scholarlypublications.universiteitleiden.nl/access/item%3A2961976/download}
\item Britannica, ``Xinca language''. \url{https://britannica.com/topic/Xinca-language}
\end{itemize}

\section{Religious and Community Resources}
I found no scripture translations or named community organisations I could verify.

\section{Digital and Audio Resources}
The Webonary site above may host lexical data. I found no audio or video resources.

\section{Curriculum Development Context}
Xinca has no first-language speakers, so a curriculum would be a revitalization effort. It would have to draw on the colonial grammar and wordlists, Sachse's research, and the work of the younger second-language learners. I found nothing on existing school programmes.
\end{document}